% Part: computability
% Chapter: recursive-functions
% Section: bounded-minimization

\documentclass[../../../include/open-logic-section]{subfiles}

\begin{document}

\olfileid[hi]{cmp}{rec}{bmi}
\olsection{परिबद्ध न्यूनीकरण}

\begin{explain}
किसी गुण अथवा संबंध~$P$ को संतुष्ट करने वाली लघुतम संख्या के रूप में
किसी फलन को परिभाषित करना प्रायः उपयोगी होता है। यदि $P$ निर्णेय है, तो
हम सभी संभव संख्याएँ $0$, $1$, $2$, \dots{} आज़माते हुए,~$P$ को
संतुष्ट करने वाली लघुतम संख्या मिलने तक, इस फलन को सरलता से संगणित कर
सकते हैं। इस प्रकार की अपरिबद्ध खोज हमें आदिम पुनरावर्ती फलनों के क्षेत्र
से बाहर ले जाती है। किंतु यदि हमारी रुचि केवल किसी \emph{स्वतंत्र रूप से
दिए गए परिबंध से छोटी} लघुतम संख्या में है, तो हम आदिम पुनरावर्ती क्षेत्र
में बने रहते हैं। दूसरे शब्दों में, और कुछ अधिक सामान्यतः, मान लें कि
हमारे पास कोई आदिम पुनरावर्ती संबंध~$R(x,z)$ है। उस फलन पर विचार करें जो
$x$ और~$y$ को उस लघुतम $z < y$ पर भेजता है जिसके लिए $R(x, z)$ है।
$R(x, 0)$, $R(x, 1)$, \dots, $R(x, y-1)$ को जाँचकर इसे भी संगणित किया
जा सकता है। किंतु यह आदिम पुनरावर्ती क्यों है?
\end{explain}

\begin{prop}
यदि $R(\vec x, z)$ आदिम पुनरावर्ती है, तो फलन $m_R(\vec{x}, y)$ भी
आदिम पुनरावर्ती है, जो ऐसा कोई मान होने पर उस लघुतम~$z$ को लौटाता है जो
~$y$ से छोटा है और जिसके लिए $R(\vec x, z)$ मान्य है, तथा अन्यथा $y$
लौटाता है। हम फलन~$m_R$ को इस प्रकार लिखेंगे:
\[
\bmin{z < y}{R(\vec{x}, z)},
\]
\end{prop}

\begin{proof}
ध्यान दें कि ऐसा कोई~$z < 0$ नहीं हो सकता जिसके लिए $R(\vec{x}, z)$
हो, क्योंकि कोई $z < 0$ होता ही नहीं। अतः $m_R(\vec x, 0) = 0$।

जब परिबंध $y + 1$ रूप का हो, तब हमारे पास तीन प्रकरण हैं:
\begin{enumerate}
\item ऐसा कोई $z < y$ है जिसके लिए $R(\vec{x}, z)$ मान्य है; इस
स्थिति में $m_R(\vec{x}, y+1) = m_R(\vec{x}, y)$।
\item ऐसा कोई~$z<y$ नहीं है, किंतु $R(\vec{x}, y)$ मान्य है; तब
$m_R(\vec{x}, y+1) = y$।
\item ऐसा कोई $z < y+1$ नहीं है जिसके लिए $R(\vec{x}, z)$ मान्य हो;
तब $m_R(\vec{z}, y+1) = y+1$।
\end{enumerate}
अतः हम $m_R(\vec x, 0)$ को आदिम पुनरावर्तन द्वारा इस प्रकार परिभाषित कर सकते हैं:
\begin{align*}
m_R(\vec{x}, 0) & = 0\\
m_R(\vec{x}, y+1) & =
\begin{cases}
m_R(\vec{x}, y) & \text{यदि $m_R(\vec x, y) \neq y$}\\
y & \text{यदि $m_R(\vec x, y) = y$ और $R(\vec{x}, y)$}\\
y+1 & \text{अन्यथा।}
\end{cases}
\end{align*}
ध्यान दें कि ऐसा $z<y$ है जिसके लिए $R(\vec x, z)$ मान्य हो, तभी और
केवल तभी जब $m_R(\vec x,
y) \neq y$।
\end{proof}

\begin{prob}
मान लें कि $R(\vec x, z)$ आदिम पुनरावर्ती है। फलन $m'_R(\vec{x}, y)$
को परिभाषित करें, जहाँ यह फलन ऐसा कोई मान होने पर उस लघुतम~$z$ को लौटाता
है जो~$y$ से छोटा है और जिसके लिए $R(\vec{x}, z)$ मान्य है, तथा अन्यथा
$0$ लौटाता है; इसे~$\Char{R}$ से आदिम पुनरावर्तन द्वारा परिभाषित करें।
\end{prob}

\end{document}
